\begin{lemma}
\label{lemma-cone-pseudo-coherent}
Let $R$ be a ring and $m \in \mathbf{Z}$.
Let $(K^\bullet, L^\bullet, M^\bullet, f, g, h)$ be a distinguished
triangle in $D(R)$.
\begin{enumerate}
\item If $K^\bullet$ is $(m + 1)$-pseudo-coherent and
$L^\bullet$ is $m$-pseudo-coherent then $M^\bullet$ is
$m$-pseudo-coherent.
\item If $K^\bullet, M^\bullet$ are $m$-pseudo-coherent, then
$L^\bullet$ is $m$-pseudo-coherent.
\item If $L^\bullet$ is $(m + 1)$-pseudo-coherent and $M^\bullet$
is $m$-pseudo-coherent, then $K^\bullet$ is $(m + 1)$-pseudo-coherent.
\end{enumerate}
\end{lemma}

\begin{proof}
Proof of (1). Choose $\alpha : P^\bullet \to K^\bullet$
with $P^\bullet$ a bounded complex of finite free modules
such that $H^i(\alpha)$ is an isomorphism for $i > m + 1$ and
surjective for $i = m + 1$. We may replace $P^\bullet$ by
$\sigma_{\geq m + 1}P^\bullet$ and hence we may assume that $P^i = 0$
for $i < m + 1$. Choose $\beta : E^\bullet \to L^\bullet$ with $E^\bullet$
a bounded complex of finite free modules such that
$H^i(\beta)$ is an isomorphism for $i > m$ and
surjective for $i = m$. By
Derived Categories,
Lemma \ref{derived-lemma-lift-map}
we can find a map $\gamma : P^\bullet \to E^\bullet$ such that the diagram
$$
\xymatrix{
K^\bullet \ar[r] & L^\bullet \\
P^\bullet \ar[u] \ar[r]^\gamma & E^\bullet \ar[u]_\beta
}
$$
is commutative in $D(R)$. The cone $C(\gamma)^\bullet$ is a bounded
complex of finite free $R$-modules, and the commutativity of the
diagram implies that there exists a morphism of distinguished triangles
$$
(P^\bullet, E^\bullet, C(\gamma)^\bullet)
\longrightarrow
(K^\bullet, L^\bullet, M^\bullet).
$$
It follows from the induced map on long exact cohomology sequences and
Homology, Lemmas \ref{homology-lemma-four-lemma} and
\ref{homology-lemma-five-lemma}
that $C(\gamma)^\bullet \to M^\bullet$ induces an isomorphism
on cohomology in degrees $> m$ and a surjection in degree $m$.
Hence $M^\bullet$ is $m$-pseudo-coherent.

\medskip\noindent
Assertions (2) and (3) follow from (1) by rotating the distinguished
triangle.
\end{proof}
